% Root and B full literal audits PASS; these cuts contain no standalone table-census assertion.
\subsection{Additional profile cuts for the four-object search}

Three elementary reductions supplement the accepted profile lemmas in
the seventeen-arrow four-object search.  Each is applied to an actual
category; none assumes that an arbitrary quotient has finite homotopy type.

\paragraph{A full singleton cut.}
Suppose the object set is partitioned into two nonempty parts $A,B$,
there are no arrows from $B$ to $A$, and every hom-set from an object
of $A$ to an object of $B$ is a singleton.  The category is exactly
the categorical join of its two full subcategories.  Hence its nerve
realization is $BA*BB$.  Choosing a vertex in each nonempty space
and collapsing the union of the two cones on those vertices gives
the usual homotopy equivalence with the suspension of the smash
product.  The collapsed subspace is contractible and is a CW
subcomplex.  Thus all products of positive-degree reduced
cohomology classes vanish, without connectedness assumptions on
the two parts.

\paragraph{The six-arrow strong pair with a singleton side.}
Suppose two objects $a,b$ have corners of order two and exactly
one arrow in each direction.  In a minimal-arrow projective-plane
candidate there is no proper retract.  Consequently the composites
of those two cross arrows are the nonidentity idempotents of their
corners: if either composite were an identity, there would be a
proper retract.  In particular each corner is $\{1,e\}$ with $e$
a zero, and has no unit other than its identity.  If every incoming
hom-set at $b$ from another object is empty or a singleton, its
nonidentity idempotent fixes all incoming cross arrows, literally.
Lemma~\ref{C:one-sided-corner-deletion} therefore deletes $b$ by an
actual homotopy equivalence.  The same conclusion holds if every
outgoing cross hom-set at $b$ is empty or a singleton, by opposition.
At seventeen arrows the result has at most thirteen arrows, since
the two corner arrows and the two internal cross arrows disappear.
It is excluded by Theorem~\ref{C:CP2-lower17}.  The alternative
that the original category already has a retract is excluded by
the same lower bound.

\paragraph{Trees with low-dimensional internal corners.}
Let the undirected support of an acyclic category be a tree.
No two nonidentity cross arrows are composable, since transitivity
would supply a third support edge and hence a triangle.  The nerve
is therefore the union of its two-object edge subcategories, glued
along the corner nerves.  This is an iterated homotopy pushout of
CW subcomplexes along the support tree.

Assume every internal support vertex has a corner monoid of order
at most five.  The accepted classification through order five
gives $H_n(BM;\mathbb Q)=0$ for every such corner and every $n\geq3$.
In the Mayer--Vietoris sequence of the edge union, the intersections
that actually glue two pieces are precisely these internal corners.
It follows for every $n\geq4$ that
\[
 H_n(BC;\mathbb Q)\cong
 \bigoplus_{\{a,b\}\text{ a support edge}}
 H_n(BC|_{\{a,b\}};\mathbb Q).
\]
In particular bounded rational homology of $BC$ passes to every
edge piece, and restriction in fourth rational cohomology is
injective.  It suffices that every edge piece have zero products
of two degree-two classes.  Each of the following numerical
conditions supplies that conclusion:
\begin{enumerate}
\item its cross hom-set has at most two elements, by the universal
two-cross-arrow theorem (the singleton case is a join);
\item both endpoint corners have order at most four, so their
second rational cohomology vanishes and every degree-two class
is relative to the vertex spaces; the relative quotient is a
suspension;
\item one endpoint corner has order at most two and the other has
order at most eleven.  The first corner has projective trivial
left and right rational modules.  Its actual orbit collapse is a
rational cohomology-ring equivalence to an ordinary act cone over
the other corner.  The cone has bounded rational homology by the
preceding comparison, and Lemma~\ref{C:eleven-bounded-all-products}
annihilates all its degree-two products.
\end{enumerate}
All three conditions are statements about every degree-two product
on the edge piece.  In particular no second-cohomology rank-one
assumption is passed from the whole tree to its pieces.
